\documentclass[]{jlreq}
% --- Engine guard for LuaLaTeX ---
\providecommand{\epTeXinputencoding}[1]{}
% --- end guard ---
\usepackage{luatexja}
\usepackage{amsmath,amssymb,amsthm}
\usepackage{tikz-cd}
\usepackage{geometry}
\usepackage{hyperref}
\geometry{margin=25mm}
\title{BTop類における初期位相・終期位相と普遍性\thanks{Lua\LaTeX\ でのコンパイルを想定しています。}}
\author{CODEX (OpenAI) -- Leanファイル生成 \and CODEX (OpenAI) -- TeXファイル生成}
\date{2025年9月28日}

\theoremstyle{definition}
\newtheorem{definition}{定義}
\theoremstyle{remark}
\newtheorem{remark}{補足}

\begin{document}
\maketitle

\begin{abstract}
学部生の読者を対象に、Bourbaki 流に定式化したトポロジー圏 \(\mathbf{BTop}\) における初期位相と終期位相を解説する。
Lean 4/mathlib を用いた実装 (`BTopCategory.lean`, `BTopInitialFinal.lean`) を参照しながら、積空間・直和空間・商空間などの普遍性を図式で示す。
\end{abstract}

\tableofcontents

\section{導入}
Bourbaki の枠組みでは、位相空間は集合とその開集合族の組として与えられ、連続写像は「写像＋連続性証明」の順序対で表現される。
Mathlib の `BTopCategory.lean` では、この考え方を Lean 上の圏として実装し、$
\mathbf{BTop}$ と通常のトポロジー圏 $\mathbf{Top}$ の同値を構成した。
本稿では、その上に構築された `BTopInitialFinal.lean` を手がかりに、初期位相と終期位相の普遍性を学生向けに整理する。

\section{BTop と Top の対応}
\begin{definition}
`BTop` は以下のデータを持つ：
\begin{itemize}
  \item オブジェクト：型 \(\alpha\) と Bourbaki 位相 \(\tau\) の対 \((\alpha,\tau)\)
  \item 射：連続写像 \(f : (\alpha,\tau) \to (\beta,\sigma)\) とその連続性証明の順序対
\end{itemize}
`BTopCategory.lean` ではこれを Lean の圏として実装し、忘却関手
\(
  \mathrm{forgetToTopCat} : \mathbf{BTop} \to \mathbf{Top}
\)
と逆方向の `ofTopCat` を構成した。
\end{definition}

\begin{center}
\begin{tikzcd}
\mathbf{BTop} \arrow[r, shift left=1.4, "\mathrm{forgetToTopCat}"] \arrow[d, "\mathrm{forget}"']
  & \mathbf{Top} \arrow[l, shift left=1.4, "\mathrm{ofTopCat}"] \arrow[d, "\Gamma"] \\
\mathbf{Set} \arrow[r, equal] & \mathbf{Set}
\end{tikzcd}
\end{center}
ここで縦方向は集合忘却関手である。同値 $\mathrm{equivTopCat} : \mathbf{BTop} \simeq \mathbf{Top}$ によって、$\mathbf{Top}$ で既知の構造（極限・余極限、随伴など）を $\mathbf{BTop}$ に輸送できる。

\section{初期位相の構成}
$X$ を集合、$F_i = (X_i,\tau_i)$ を $\mathbf{BTop}$ の対象、$\varphi_i : X \to X_i$ を写像の族とする。

\begin{definition}[初期位相]
$\varphi_i$ をすべて連続にする最弱の位相を \emph{初期位相} と呼び、Lean では
\[
\texttt{initialTopology F φ} = \texttt{ofTopologicalSpace} \Bigl(\bigwedge_{i\in I} \mathrm{TopologicalSpace.induced}(\varphi_i, \tau_i)\Bigr)
\]
と定義する。これを持つ $X$ を `initial F φ` と書く。
\end{definition}

\begin{remark}[普遍性]
任意の $Z\in\mathbf{BTop}$ と写像 $f : Z \to X$ に対し、
\[
 f \text{ が初期位相に関して連続} \iff \forall i,~\varphi_i\circ f \text{ が } F_i \text{ へ連続}
\]
が成り立つ。Lean では `initial_continuous_iff` が対応する。
\end{remark}

\section{終期位相の構成}
次に、写像 $\psi_i : X_i \to X$ の族を考える。

\begin{definition}[終期位相]
全ての $\psi_i$ を連続にする最強の位相を \emph{終期位相} と呼び、Lean では
\[
\texttt{finalTopology F ψ} = \texttt{ofTopologicalSpace} \Bigl(\bigvee_{i\in I} \mathrm{TopologicalSpace.coinduced}(\psi_i, \tau_i)\Bigr)
\]
と定義し、$X$ に付けたものを `final F ψ` と書く。
\end{definition}

\begin{remark}[普遍性]
任意の $Z\in\mathbf{BTop}$ と写像 $f : X \to Z$ に対し、
\[
 f \text{ が終期位相に関して連続} \iff \forall i,~f\circ\psi_i \text{ が } F_i \text{ から } Z \text{ へ連続}
\]
が成り立つ。Lean では `final_continuous_iff` が実装されている。
\end{remark}

\section{積と初期位相：\texttt{piObj}}
`BTopInitialFinal.lean` では家族 $F : I \to \mathbf{BTop}$ に対して
\begin{align*}
  \texttt{piObj}~F &= \texttt{initial}\bigl(F, \; \varphi_i : (\prod_j X_j) \to X_i\bigr),\\
  \texttt{piProj}~F~i &: \piObj F \to F_i,
\end{align*}
を定義している。ここで $\varphi_i$ は第 $i$ 成分を取り出す写像である。

\begin{center}
\begin{tikzcd}
& Z \arrow[dl, "f"'] \arrow[drr, dashed, "f^{\\sharp}" description] & \\
X_i &&& \piObj F \arrow[lll, "\pi_i" ]
\end{tikzcd}
\end{center}

\noindent Lean の補題 `continuous_to_pi_iff` は、図式が可換になる条件が各成分の連続性に帰着することを保証する。
\[
  f : Z \to \piObj F \text{ が連続 }\iff \forall i,~(\pi_i \circ f) : Z \to F_i \text{ が連続}
\]
これにより、$\piObj F$ が $\mathbf{BTop}$ における積の普遍性を満たすことがわかる。

\section{直和と終期位相：\texttt{sigmaObj}}
直和空間は終期位相を使って
\[
  \texttt{sigmaObj}~F = \texttt{final}\bigl(F, \; \iota_i : F_i \to \Sigma_j X_j\bigr)
\]
で定義され、各包含写像を `sigmaInl` として Lean に実装している。

\begin{center}
\begin{tikzcd}
F_i \arrow[r, "\iota_i"] \arrow[drr, bend right=20, "g_i"'] & \sigmaObj F \arrow[dr, dashed, "g" ] & \\
&& Z
\end{tikzcd}
\end{center}

`continuous_from_sigma_iff` により、$g : \sigmaObj F \to Z$ が連続であることは、各成分 $g_i = g \circ \iota_i$ の連続性に同値である。したがって $\sigmaObj F$ は余積（コプロダクト）の普遍性を満たす。

\section{商位相：\texttt{quotientObj}}
$X\in\mathbf{BTop}$ と同値関係（集合論の意味での）$\sim$ による集合商 $\mathrm{Quot}(X,\sim)$ を考える。
Lean では終期位相を用いて
\[
  \texttt{quotientObj}~X~r = \texttt{final}\bigl(\{X\}, \; q : X \to \mathrm{Quot}(r)\bigr)
\]
を与え、射 `quotientMap : X \to quotientObj X r` が連続になるよう定義している。

\begin{center}
\begin{tikzcd}
X \arrow[r, "q"] \arrow[d, "f"'] & \mathrm{Quot}(X,r) \arrow[dl, dashed, "\overline{f}"] \\
Z &
\end{tikzcd}
\end{center}

補題 `continuous_from_quotient_iff` は、$\overline{f}$ が連続であることが $f$ の連続性と同値であることを示す。これが商位相の普遍性そのものである。

\section{Lean 実装との対応表}
\begin{center}
\begin{tabular}{ll}
\hline
数学的概念 & Lean 中の識別子 \\
\hline
初期位相 & \texttt{initialTopology}, \texttt{initial} \\
終期位相 & \texttt{finalTopology}, \texttt{final} \\
積空間 & \texttt{piObj}, \texttt{piProj}, \texttt{continuous\_to\_pi\_iff} \\
直和空間 & \texttt{sigmaObj}, \texttt{sigmaInl}, \texttt{continuous\_from\_sigma\_iff} \\
商空間 & \texttt{quotientObj}, \texttt{quotientMap}, \texttt{continuous\_from\_quotient\_iff} \\
随伴 & \texttt{disc\_adj\_forget}, \texttt{forget\_adj\_indisc} \\
\hline
\end{tabular}
\end{center}

\section{演習問題}
\begin{enumerate}
  \item $F_1,F_2\in\mathbf{BTop}$ について、$\piObj$ の普遍性から積位相の基本開集合系を記述してみよ。
  \item 離散位相を持つ集合族から構成した `sigmaObj` が、集合としての直和に離散位相を貼り付けたものと一致することを示せ。
  \item `quotientObj` を用いて、円周 $S^1$ を区間 $[0,1]$ の端点貼り合わせとして実装するアイデアを考察せよ。
\end{enumerate}

\section{参考文献}
\begin{itemize}
  \item Nicolas Bourbaki, \textit{\'El\'ements de math\'ematique}, Topologie g\'en\'erale.
  \item Lean/mathlib documentation: \url{https://leanprover-community.github.io/mathlib4_docs/}
\end{itemize}

\end{document}
